\documentclass[11pt]{article}
\usepackage{amsmath,amssymb,amsthm,mathtools,hyperref}

\begin{document}
\section*{Mechanisms of parasite manipulation of host behaviour}
\label{prob:330}
\label{rank:210}
\par\smallskip

\noindent\textbf{Domain.} Host--parasite interactions; causal behaviour attribution. \quad \textbf{Retained status/source context.} Open system-specific causal attribution. Documented behavioural changes and tissue localization do not establish a complete manipulation pathway or its adaptive benefit.

\paragraph{Definitions and assumptions.} For a declared host--parasite system let $A$ be natural infection/exposure status, $M$ a predeclared standardized host neural, endocrine or muscle-state summary, $B$ a declared scalar behaviour score or event indicator and $F$ parasite reproductive-success outcome over horizon $H$. Baseline host state $C$, infection burden and sickness-related changes are observed or assigned explicit latent-variable models. A restricted structural model specifies $P(M,B,F\mid A,C)$ and measurement error. Its causal specification must also state intervention semantics, temporal ordering and restrictions on latent variables; the observational conditional distribution alone does not determine counterfactual laws. Observing burden or sickness does not by itself justify adjustment for them. Behavioural manipulation means a causal behaviour change with a declared parasite-benefit criterion; it is distinct from generic host impairment. Only already documented coarse variables are considered.

\paragraph{Problem statement.} Determine which causal claims are identifiable from existing evidence:
    \[\Delta_B=\mathbb{E}[B^{A=1}-B^{A=0}],\qquad
    \Delta_F(b,b')=\mathbb{E}[F^{A=1,B=b}-F^{A=1,B=b'}].\]
    The superscripts denote specified counterfactual exposures or behavioural states under consistency and justified confounding assumptions. When identification assumptions fail but the counterfactual contrasts remain well-defined, quantify bounds over the causal models compatible with the declared evidence and retained assumptions. Specify any outcome-support or sensitivity restrictions used, and distinguish population identification bounds from finite-sample uncertainty. Failure to define intervention versions or consistent counterfactual outcomes is instead a failure to define the estimand, not merely a reason for wider bounds. A mediator pathway requires additional identification beyond the infection-to-behaviour and behaviour-to-parasite-success contrasts.

\paragraph{Additional binding assumptions and resolution criterion.} Give behavior a declared scalar score or indicator before subtracting its potential outcomes. Evaluate the behavior-to-parasite-success contrast within a fixed infected context, with the same burden and baseline conventions, rather than mixing infected and uninfected hosts. The averaging population and the meaning of the shared burden convention remain to be declared for an instance: setting $A=1$ alone does not specify selection of naturally infected hosts, and a common burden convention does not by itself select standardization, conditioning or intervention. Declare burden timing and whether later burden can respond to behavior. Specify whether $A$ denotes exposure or established infection, its intervention versions and timing, the implementations and admissible values of $b,b'$, reproductive-success units, the horizon $H$ and its origin, and conventions for host death, parasite loss, censoring and missing behavior. Declare the expectation domain under the stated joint causal law; a scalar score alone does not ensure that an expected difference exists. The infection-to-behaviour and behaviour-to-parasite-success contrasts do not identify a mediated causal path or historical adaptation. Explicitly compare generic sickness, host-beneficial responses and parasite-benefit explanations within the observed model class. For an attribution of the infection-associated behavior change, explain how the chosen behavioral comparison bears on that change and apply the declared parasite-benefit criterion; benefit for an unrelated pair of behavioral values is not sufficient attribution. The competing explanations may coexist or remain observationally indistinguishable.

\paragraph{Alternative formulations.} (i) A host-pathology model predicts behaviour from generalized illness without a parasite-benefit claim. (ii) A parasite-adaptation model additionally predicts a supported fitness contrast under behaviour changes. (iii) A mechanistic mediation model distinguishes measured host-state pathways, while a comparative phylogenetic model tests repeated associations. These analyses answer different causal and historical questions.

\paragraph{Obstructions and partial progress.} Three-dimensional observations in manipulated ants reveal extensive fungal association with muscles and no detected fungal cells inside the examined brains. Thus direct brain invasion is not a required default explanation; absence of detected invasion does not establish absence of biochemical or functional neural effects. Localization alone does not identify the biochemical or physiological causal chain, and host distress can mimic targeted behavioural effects. Parasite transmission opportunities and host ecology also confound inferred benefit. The remaining question is defensible attribution and pathway discrimination for a specified natural system, with adaptation claims separated from pathology and correlation.

\paragraph{Sources (supplied provenance; not all rechecked).} \begin{itemize}
\item Supplied Wikipedia background: behaviour-altering parasite. \url{https://en.wikipedia.org/wiki/Behavior-altering_parasite} (Retained for provenance; not used as scientific verification of current open status.)
\item Dobson, Holt and Tilman (editors), Unsolved Problems in Ecology (2020). \url{https://www.jstor.org/stable/j.ctvs9fh2n} (Book metadata and contents checked; chapter text access restricted. Broad background, not this exact mathematical statement.)
\item Fredericksen et al., Three-dimensional visualization and a deep-learning model reveal complex fungal parasite networks in behaviorally manipulated ants (2017). \url{https://pubmed.ncbi.nlm.nih.gov/29114054/} (Primary tissue-localization study checked; spatial observations do not establish a complete causal pathway.)
\end{itemize}

\section*{Binding catalogue conventions}
This edition revises both supplied catalogues and preserves original problem numbers 1--330. It adds one globally unique label \verb|\label{rank:x}| per problem, where smaller $x$ means greater provisional difficulty. A rank label is an anchor at that problem's section: \verb|\ref{rank:x}| returns the section number, not the literal difficulty number. The visible \texttt{rank:x} text and the index give the difficulty number itself.

The entries are \emph{parameterized mathematical research schemes}, including empirical prediction, causal identification, historical reconstruction, theory and engineering. They are not 330 certified unsolved mathematical conjectures. A rigorous instance fixes the objects and restrictions below; changing those objects changes the question. Established identities, finite calculations, non-identification constructions and solved special cases are not being posed anew as open. Some residual tasks concern informative bounds, biological realization or external validation rather than a new theorem.

Every main definition and problem statement received a mathematical editorial review. Each entry now has an additional binding-conditions paragraph; direct corrections are separately recorded in the accompanying audit. The alternative formulations inherit the applicable conventions, but are different targets unless equivalence is proved. Where an old broad phrase conflicts with an explicit restriction below it, the explicit restriction controls.

\textbf{Literature boundary.} Thirteen selective primary-literature checks cover eighteen entries in this revision; their access levels and scope are recorded explicitly. The supplied bibliography and its earlier source-access descriptions are retained as provenance, not represented as newly verified. Neither this review nor a source's description proves the absence of later solutions. The date of this edition is 13 October 2026; there is no claim of exhaustive current-status certification for all 330 entries.

\subsection*{Binding conventions for a concrete instance}
\paragraph{Objects, quantifiers and non-vacuity.}
An instance declares measurable state, input, output and parameter spaces; a fixed candidate family; admissible initial laws, interventions and environments; a finite horizon or specified asymptotic regime; observation and sampling laws; and performance metrics, positive normalization scales and positive tolerances. These are chosen before fitting or testing, not adjusted after seeing failures. Shared parameters must genuinely be shared across the declared validation domain. Finite dimension alone does not restrict arbitrary real-number encodings: specify a finite library, bounded description length, or regularity/complexity bounds. Exclude constant or no-signal solutions through a fixed informative-output or baseline-improvement requirement where appropriate. A universally quantified result cannot be established by choosing one trivial favorable instance.

\paragraph{Probability laws and data.}
Every expectation or probability refers to the specified stochastic law, including initial uncertainty, process noise, observation noise and sampling as appropriate. A deterministic equation is not a probability model without a declared random input or initial law. ``Observed'' targets in population-level formulas mean population observables $g_j^*$, not noisy estimates asserted to be exact. Let $\widetilde g_j$ be independent validation estimates and suppose
\[
\mathbb{P}\{\lvert\widetilde g_j-g_j^*\rvert\le r_j\text{ for every }j\}\ge1-\alpha.
\]
Then an uncertainty-aware sufficient certification rule for positive target tolerances $\varepsilon_j$ is
\[
\max_j\frac{|\widehat g_j-\widetilde g_j|+r_j}{\varepsilon_j}\le1.
\]
On that simultaneous confidence event it implies every population error is at most its tolerance, by the triangle inequality. Vector/norm versions use corresponding valid simultaneous bounds. Training uncertainty must also be included when the claim averages over repeated training data. Evaluation on finitely many conditions does not establish a continuous-domain supremum without covering and regularity bounds.

\paragraph{Dynamics and well-posedness.}
For ODEs and SDEs give coefficient domains, initial laws and existence/uniqueness conditions in the intended solution sense, or keep the entire admitted solution set. For stochastic jump models specify nonnegative off-diagonal rates, diagonal normalization, rate orientation and non-explosion. For PDEs specify spatial domains, boundary/initial conditions, conservation or source terms, and a weak or classical solution class. Moving-domain concentration equations include dilution/transport terms. A stationary law is used only for a time-homogeneous augmented model under existence assumptions; a stationary approximation to a driven system must be justified. Absorbing extinction states are retained. Survival-conditioned, quasi-stationary and unconditional laws are distinct and must not be substituted for one another.

\paragraph{Identification and causal meaning.}
Two parameter values are observationally equivalent when they induce the same full distribution under all allowed observation/intervention designs, not merely the same likelihood value for one sample. A target is identified only if constant on the equivalence class. Gauge symmetries, label switching and reference conventions are quotiented or fixed. Finite-data compatible sets use a declared confidence, likelihood or discrepancy rule; an empty set signals model rejection. Bayesian credible sets are not automatically confidence sets. A prior or regularizer does not create data identification. Bounds intended to resolve a question must be informative at a predeclared width or decision threshold.

A causal contrast uses well-defined intervention versions and a stated causal model. Random assignment, or explicit consistency, overlap and exchangeability assumptions, is required for causal interpretation from observations. Interference and time-varying confounding are represented when present. Fixing a mediator by intervention differs from conditioning on its observed value. Marginal potential-outcome laws do not determine individual cross-world effects or an arbitrary joint coupling. Ethically admissible designs, rather than unconstrained interventions, define the experimental menu.

\paragraph{Metrics, units, optimization and limits.}
Specify the sampling or reference measure for densities and integrals, the log base, and dimensionless arguments of logarithms. Total variation is $\sup_A|P(A)-Q(A)|$; Wasserstein distances require the declared ground metric and appropriate moment conditions; KL divergence uses its stated absolute-continuity convention. All normed readouts are scaled consistently, and error probabilities use $0<\alpha<1$. A ratio has a positive denominator or an explicit undefined/zero convention. Thermodynamic entropy production is nonnegative in expectation under the relevant hypotheses; individual stochastic realizations need not be. An infimum is not an attained minimum without compactness/coercivity and continuity or another existence argument. Fix the order of large-population, long-time and vanishing-mutation limits. Stability claims declare perturbation space and symmetries.

\subsection*{How the provisional difficulty order was assigned}
Four editorial components are scored from 1 (lower burden) to 5 (higher burden): (a) formal or inferential difficulty, (b) access to evidence and identification burden, (c) mechanistic/computational integration, and (d) independent validation and transfer. Their unweighted total orders the entries. For selected entries in the highest bands, an explicit editorial priority breaks equal totals; all other ties are broken by original problem number. Equal totals are \emph{not} evidence of reliably different true difficulties. Scores are judgments, not measurements, completion-time estimates or consensus rankings.

The object ranked is an informative, nondegenerate resolution of the \emph{scoped} scientific target, with the stated budget and validation domain. A sharp useful impossibility result can qualify; an already known identity, vacuous bound, arbitrary favorable model or lookup table cannot. Rank 330 is merely the least difficult relative assessment in this catalogue, not a claim of an easy biological solution. Broad historical titles do not automatically outrank precisely formulated mechanistic tasks.



\end{document}
